% META: {"id":"TEXP-ARI-CO-002","chapitre":"TEXP-ARITHMETIQUE","type_objet":"corrige","exercice_ref":"TEXP-ARI-EX-002","capacites_codes":["C2"],"sources_inspiration":[],"status":"approved","origin":"generated","status_history":["generated","audited","accepted","approved"],"release_acceptance":"RELEASE_OWNER_BATCH_ACCEPTANCE","acceptance_closure_digest":"sha256:9b3ccf9a81c5520fbb7e03b7b2d3e7bf2057a02834b6d908bf3be5f49f3b553f","acceptance_receipt_digest":"sha256:4fad86f0282e0fa4e0419cca1ad6379ae7af5a280c04a4a90880f90a1c044242"}
% BEGIN-VERIFY
% assert (4*3) % 11 == 1
% solutions = [x for x in range(11) if (4*x) % 11 == 7 % 11]
% assert solutions == [10]
% END-VERIFY
\begin{corrige}{TEXP-ARI-EX-002}
\textbf{1.} $4\times3=12\equiv1\ [11]$ : $3$ est bien l'inverse de $4$ modulo $11$.

\textbf{2.} En multipliant les deux membres de $4x\equiv7\ [11]$ par $3$ : $12x\equiv21\ [11]$, soit $x\equiv21\ [11]$ (car $12\equiv1\ [11]$), c'est-à-dire $x\equiv10\ [11]$.
\end{corrige}
